\documentclass[11pt,letterpaper]{article}
\usepackage[utf8]{inputenc}
\usepackage[T1]{fontenc}
\usepackage[spanish,es-tabla]{babel}
\usepackage[margin=2.4cm,headheight=15pt]{geometry}
\usepackage{lmodern,microtype,amsmath}
\usepackage{graphicx,booktabs,longtable,tabularx,array}
\usepackage{xcolor,listings,fancyhdr}
\usepackage[hidelinks]{hyperref}
\usepackage{url}
\urlstyle{same}
\setlength{\parindent}{0pt}
\setlength{\parskip}{0.35em}
\setlength{\emergencystretch}{4em}
\renewcommand{\arraystretch}{1.12}
\setcounter{tocdepth}{2}
\setcounter{secnumdepth}{2}
\newcommand{\artifactdir}{generado-ampliado-20261102T050747Z-2b7e3c41}

\pagestyle{fancy}
\fancyhf{}
\fancyhead[L]{\small MongoDB y PostgreSQL JSONB}
\fancyhead[R]{\small Informe parcial}
\fancyfoot[C]{\thepage}
\hypersetup{
 pdftitle={MongoDB y PostgreSQL JSONB: informe académico parcial},
 pdfauthor={Jesel Moreno Zuniga}
}
\newcommand{\pendiente}[1]{%
 \par\smallskip\noindent
 \fbox{\parbox{0.94\linewidth}{\small\textbf{Pendiente:} #1}}
 \par\smallskip}
\lstset{
 basicstyle=\ttfamily\scriptsize,
 breaklines=true,columns=fullflexible,keepspaces=true,
 showstringspaces=false,frame=single,
 aboveskip=0.6em,belowskip=0.6em
}
\newcommand{\reportfigure}[3]{%
 \IfFileExists{\artifactdir/#1}{%
 \begin{figure}[!htbp]\centering
 \includegraphics[width=0.92\linewidth]{\artifactdir/#1}
 \caption{#2}\label{#3}\end{figure}
 }{\pendiente{Figura no generada. #2}}}
\begin{document}
\begin{titlepage}\centering
{\Large Universidad Politécnica de Chiapas\par}
\vspace{0.6cm}
{\large Licenciatura en Ingeniería en Tecnologías de la Información
e Innovación Digital\par}
\vspace{0.7cm}
{\large Base de datos en la nube\par}
\vspace{0.9cm}
{\LARGE\bfseries Análisis comparativo del rendimiento de MongoDB
y PostgreSQL JSONB bajo restricciones de CPU y memoria
en agregaciones de documentos JSON anidados\par}
\vspace{0.7cm}
{\Large Informe académico parcial\par}
\vspace{0.7cm}
\begin{tabularx}{0.92\linewidth}{rX}
Estudiante: & Jesel Moreno Zuniga\\
Grupo: & 7B\\
Profesor: & Horacio Irán Solís Cisneros\\
Actividad: & 2026.C3.EVA02
\end{tabularx}
\vfill
{1 de noviembre de 2026\par}
\smallskip
\begin{minipage}{0.91\linewidth}\small
La matriz prevista comprende 108 ejecuciones. El responsable confirma
dos mediciones definitivas elegibles. El documento distingue la evidencia
cotejada localmente, la información comunicada y los pendientes.
\end{minipage}
\end{titlepage}

\section*{Resumen}
Se estudia el comportamiento de MongoDB y PostgreSQL JSONB en
agregaciones sobre 500\,000 documentos JSON anidados. El diseño incluye
tres perfiles de CPU y memoria, dos consultas, tres concurrencias y
tres repeticiones, para 108 ejecuciones previstas en la matriz base.
El protocolo establece calentamiento de 60 segundos, medición de
300 segundos y enfriamiento de 60 segundos.

El responsable confirma dos mediciones definitivas elegibles para Q1,
perfil A, concurrencia 1 y repetición 1. Las tablas documentan su cotejo
contra CSV y archivos de commit existentes. Se registran throughput,
percentiles, CPU throttling, RSS máximo, spill y estado de ejecución,
sin completar datos ausentes con cero.

El alcance es descriptivo: una observación por motor no permite inferir
variación entre repeticiones, efectos de concurrencia o superioridad
general. Se conserva la limitación del índice textual de PostgreSQL
frente al cast temporal de Q2. También se documenta que la reutilización
del preflight se bloqueó al cambiar de motor, sin atribuir un ahorro
de tiempo no demostrado.

\textbf{Palabras clave:} MongoDB, JSONB, agregaciones, cgroups v2,
latencia, RSS, throttling, reproducibilidad.
\tableofcontents
\clearpage

\section*{Introducción y justificación}
\addcontentsline{toc}{section}{Introducción y justificación}
El soporte JSON en sistemas relacionales y las bases orientadas a
documentos permiten formular agregaciones sobre estructuras similares.
Sin embargo, la equivalencia funcional no implica igualdad de planes,
representación física o administración de memoria. Una comparación
defendible requiere conservar la fuente de datos, la semántica y las
condiciones que explican cada observación.

Este trabajo examina dos cargas: desanidado y agrupamiento de ítems,
y filtrado temporal con ordenamiento y Top-K. La restricción explícita
de CPU y memoria permite plantear preguntas sobre capacidad, latencia
y trabajo temporal, sin asumir que una sola cifra representa todo el
comportamiento de un motor.

La entrega es parcial. Las operaciones ejecutadas dentro de una ventana
no son repeticiones independientes del escenario. Los pilotos previos,
las mediciones definitivas y los escenarios pendientes deben mantenerse
separados. Se distinguen hechos comunicados por el responsable,
archivos examinados durante la generación y limitaciones pendientes.

\textbf{Objetivo general.} Comparar MongoDB y PostgreSQL JSONB mediante
agregaciones sobre 500\,000 documentos anidados bajo cuotas de CPU y
memoria controladas.

\textbf{Objetivos específicos.} Verificar equivalencia de datos;
conservar consultas originales; medir rendimiento y señales de recursos;
describir índices y planes existentes; y delimitar el alcance de las
conclusiones de una matriz incompleta.

\input{\artifactdir/01-entorno.tex}
\input{\artifactdir/02-metodologia.tex}
\input{\artifactdir/03-resultados.tex}
\input{\artifactdir/04-graficas.tex}
\input{\artifactdir/05-planes.tex}
\input{\artifactdir/06-interpretacion.tex}
\input{\artifactdir/07-recomendaciones.tex}
\input{\artifactdir/08-conclusiones.tex}
\clearpage
\addcontentsline{toc}{section}{Referencias}
\begin{thebibliography}{9}
\bibitem{pgmemory} PostgreSQL Global Development Group. \textit{PostgreSQL 18: Resource Consumption}. \url{https://www.postgresql.org/docs/18/runtime-config-resource.html}.
\bibitem{pgexplain} PostgreSQL Global Development Group. \textit{PostgreSQL 18: Using EXPLAIN}. \url{https://www.postgresql.org/docs/18/using-explain.html}.
\bibitem{pgindex} PostgreSQL Global Development Group. \textit{PostgreSQL 18: Indexes on Expressions}. \url{https://www.postgresql.org/docs/18/indexes-expressional.html}.
\bibitem{wiredtiger} MongoDB. \textit{WiredTiger Storage Engine}. \url{https://www.mongodb.com/docs/manual/core/wiredtiger/}.
\bibitem{mongolimits} MongoDB. \textit{Aggregation Pipeline Limits}. \url{https://www.mongodb.com/docs/manual/core/aggregation-pipeline-limits/}.
\bibitem{docker} Docker. \textit{Resource constraints}. \url{https://docs.docker.com/engine/containers/resource_constraints/}.
\bibitem{kernel} The Linux Kernel documentation. \textit{Control Group v2}. \url{https://docs.kernel.org/admin-guide/cgroup-v2.html}.
\end{thebibliography}
\appendix
\input{\artifactdir/09-anexos.tex}
\end{document}
